% RESEARCH_DRAFT_SHA256: af672efd9fa8332621c0896eaff57c0ce12cba7311ea1b36d28f12fb2e2ed798
% Layout and source-reference wrappers expanded for the web renderer.
\documentclass[11pt]{article}
\usepackage[margin=25mm]{geometry}
\usepackage{fontspec,unicode-math}
\setmainfont{Latin Modern Roman}
\setmathfont{Latin Modern Math}
\usepackage{amsmath,amssymb,mathtools,xurl,hyperref}
\hypersetup{colorlinks=true,urlcolor=blue,linkcolor=blue}
\setcounter{secnumdepth}{0}
\setlength{\parindent}{0pt}
\setlength{\parskip}{5pt}
\setlength{\emergencystretch}{3em}
\newcommand{\sref}[2]{\hyperlink{source:#1:#2}{[S#2]}}
\newcommand{\cataloguescope}[1]{\paragraph{Source record.}#1}
\newcommand{\floor}[1]{\left\lfloor#1\right\rfloor}
\begin{document}
\section*{Research-attempt provenance}
This standalone research attempt retains the complete original definition body below,
including its historical source and preparation notes. Those notes describe the
earlier catalogue preparation; the new attempt and its verification are identified
separately after that body. No original ID, catalogue status, or source statement is changed.

\section{Low-Rank Subexponential Electronic-Structure Algorithms}\label{statement}

\subsection{Definitions and mathematical statement}

\paragraph{Shared notation.}Standard number systems are $\mathbb N,\mathbb Z,\mathbb Q,\mathbb R,\mathbb C$. Specialized notation and index ranges follow the cited source and the definitions below.


The electronic Schrodinger equation is an eigenvalue problem $H\psi=E\psi$ for the specified many-electron Hamiltonian. Low-rank wavefunction approximations restrict the complexity of a tensor representation of $\psi$. Subexponential runtime grows slower than $c^N$ for every fixed $c>1$ in the stated size parameter.
\paragraph{Record scope and qualifications.}
The catalogue title is retained as metadata. Named certificates, contracts, bounds or frameworks in that title are not asserted as results of the original source extract.
The Hamiltonian class, approximation norm, required precision, and size parameter are not specified in the supplied question.

Source-record limitation: The Hamiltonian class, approximation precision/norm, and input-size parameter of the requested subexponential method are absent.
\paragraph{Statement recorded in the supplied source.}
Can low-rank approximations of wavefunction parameters yield a numerical method for solving the electronic Schrödinger equation with subexponential complexity? \href{https://doi.org/10.4171/owr/2011/32}{[S2]}
\paragraph{Source status.}
The supplied catalogue labels this source record open; that metadata was checked against the supplied snapshot on 2026-10-11. The label is preserved for provenance and is not an independent verification of present open status or mathematical correctness. The catalogue literature assessment is dated 2026-08-21.
\subsection{Short English statement}
Use low-rank wavefunction representations to solve the cited electronic-structure problem in subexponential time.
\subsection{Sources}
\begin{itemize}
\item[S1] Ulam AI and the UnsolvedMath Contributors, UnsolvedMath: A Curated Collection of Open Mathematics Problems, catalogue entry. Catalogue curation and metadata: CC BY 4.0. The source extract is retained; definitions, English summaries and qualifications are editorial additions.
\url{https://huggingface.co/datasets/ulamai/UnsolvedMath}.
\item[S2] Mathematical Methods in Quantum Chemistry (2011), Oberwolfach Reports, DOI 10.4171/owr/2011/32, pp. 1769--1843
\url{https://doi.org/10.4171/owr/2011/32}.
\item Source-specific Oberwolfach Reports licences, including CC BY-SA 4.0 where applicable, are retained; see THIRD\_PARTY\_NOTICES.md.
\end{itemize}
\paragraph{Preparation and review.}
The original source and UnsolvedMath contributors retain attribution for the source extract. Codex (GPT-6; exact serving revision not exposed) assisted with definitions, English summaries, scope notes and formatting. No new human mathematical authorship or independently verified solution is asserted. Catalogue identity checked against the complete supplied registry and current repository. Distinct upstream records, including repeated questions and broad or damaged extracts, are retained for complete coverage of the source CS category. A repeated record is not a new independent conjecture. Complete mathematical-equivalence and current-literature review has not been performed. Preparation and final review follow contribution prompts C and D, with source gaps disclosed. Structural and rendering checks do not establish a new theorem or open status.
\label{end:statement}

\section{Literature and context check}
The primary report, printed pages 1810--1811, motivates low-rank wavefunction parameters but does not prescribe a tensor format, orbital order or uniform approximation norm. The following supplies a precise rank obstruction in an auxiliary finite two-state tensor model and then shows why it is not a complexity lower bound for the electronic Hamiltonian class. The distinction between a chosen tensor representation and the physical problem is essential.
\section{Research attempt: A tensor-rank obstruction and the importance of representation}

\paragraph{An exact rank requirement across a fixed cut}
For an integer $m\geq1$, on $2m$ two-state factors consider
\[
 \Psi_m=2^{-m/2}\sum_{b\in\{0,1\}^m}|b\rangle\otimes|b\rangle.
\]
Its coefficient matrix across the first-$m$/last-$m$ cut is $M=I_D/\sqrt D$, where $D=2^m$, and its squared norm is one. For an integer $0\leq r\leq D$, any approximation whose matrix across this cut has rank at most $r$ satisfies
\[
 \lVert M-B\rVert_F^2\geq1-r/D.
 \tag{1}
\]
Indeed $\ker B$ has dimension at least $D-r$. Choose that many orthonormal kernel vectors and extend to an orthonormal basis. The Frobenius norm squared is the sum of squared images on this basis; on each kernel vector $(M-B)u=Mu$ has squared norm $1/D$. This proves (1), with equality achieved by retaining any $r$ diagonal entries of $M$. Thus norm error at most $\varepsilon<1$ requires $r\geq(1-\varepsilon^2)2^m$.

Any tensor train whose central bond has dimension $r$ yields a matrix of rank at most $r$ across this cut: contracting all tensors to the left and right expresses the matrix as a sum of $r$ column-times-row products. Consequently this ordering forces an exponential central bond for fixed error. The argument does not assume a particular optimization algorithm for finding the approximation.

\paragraph{The state has a compact, exactly solvable Hamiltonian}
Let $\Phi=(|00\rangle+|11\rangle)/\sqrt2$ and define
\[
 H=\sum_{i=1}^m\left(I-|\Phi\rangle\langle\Phi|\right)_{i,m+i}.
\]
The terms act on disjoint pairs, are commuting nonnegative projectors, and each pair has a one-dimensional zero eigenspace. In a tensor product of their eigenbases, every energy is the number of pairs outside that zero eigenspace. Thus $\Psi_m$ is the unique ground state and the first excited energy is one. The Hamiltonian has $m$ explicitly specified two-factor terms, and its ground energy and factorized pair state are known without exponential computation.

If the tensor factors are reordered as $(1,m+1),(2,m+2),\ldots$, then every cut either separates whole pairs, where rank is one, or cuts one pair, where rank is two. A tensor train of bond dimension at most two follows directly by writing each Bell pair as its two-term expansion and multiplying the independent pair factors. Therefore the exponential rank obstruction depends on the chosen ordering and cut, not on an intrinsic computational difficulty of this example.

\paragraph{What this does and does not say about the source route}
This auxiliary Hamiltonian is not asserted to be a fermionic Coulomb electronic Hamiltonian in the source's class. Its role is to disprove a blanket inference from a short two-body Hamiltonian description to small tensor rank in every fixed ordering. Conversely its exact solvability disproves the inference from exponential rank in one format to exponential runtime for every method. A general electronic-structure result would need a rank bound in a specified admissible representation, a way to find that representation, controlled contraction costs, and a relation between tensor error and the requested energy or wavefunction accuracy. None follows from the word ``low-rank'' alone.

\section{Verification and adversarial audit}
The rank lower bound allows unnormalized approximants, so it also applies to normalized ones without an extra assumption. The Hamiltonian gap is obtained from disjoint commuting pair projectors, not a generic interacting-system gap theorem. Reordering gives an explicit alternative small representation. No electronic-system counterexample or universal complexity lower bound is claimed from this distinguishable-factor model.
\paragraph{Audit corrections.} Separated tensor rank, tensor ordering, algorithmic hardness and membership in the electronic Hamiltonian class.
\paragraph{External computation.} None; symbolic arguments and boundary checks only.
\section{Scope and remaining gaps}
The original electronic low-rank subexponential agenda remains unresolved. The fixed-cut obstruction and its efficient reordered representation rigorously identify missing representation hypotheses; they neither prove nor refute the intended result for the unspecified electronic class.
\paragraph{Integration review.} Codex (GPT-6; exact serving revision not exposed) reviewed the written arguments and expanded layout and source-reference wrappers. Repository metadata and the production web renderer were checked using Python 3.13 and JavaScript with MathJax 3.2.2. This model-assisted review is not independent human proof verification; no PDF compilation or formal proof checking is asserted.

\section{Completion Estimate}
\noindent\textbf{COMPLETION: 0\%}
The calculations below the preserved source statement investigate a concrete route,
but establish no substantiated new progress toward the entire original target.
Known special cases and failed sufficient conditions are not credited as a solution.
\section{References}
The original source references and catalogue attribution are preserved in the
Sources section above. The following URLs were actually consulted for this attempt;
retention of an additional original reference does not claim fresh consultation of it.
\begin{itemize}
\item \url{https://ems.press/content/serial-article-files/46348} (consulted 2026-10-11). 
\end{itemize}
\section{Model and contributor disclosure}
The parent Codex GPT-6 model has reviewed the written mathematical arguments at their declared scopes. This is model-assisted review, not independent human verification.
OpenAI gpt-6-astra, requested reasoning configuration ultra, generated this research
attempt on 2026-10-11. The exact serving revision is not exposed. Codex GPT-6
performed orchestration and repository-format review; its exact serving revision
is also not exposed. The mathematical route, boundary checks and final adversarial
self-audit followed contribution prompts A then B in
\texttt{docs/contribution\_prompts.md}, after reading \texttt{CONTRIBUTING.md}.
Python 3.14.2 assembled the standalone documents; any mathematical computation is
specified in Verification. No theorem prover was used. No human mathematical
authorship, substantial human proof correction, or independent human proof checking
is asserted. Model-assisted review, source retrieval, and structural checks do not
establish correctness or novelty. Full private reasoning transcripts are unavailable
as reproducibility artifacts; the complete arguments, inputs and outcomes actually
used are stated here or in the companion batch audit.
\section{Final status}
Unresolved as a new research contribution. This is a concrete route investigation,
not a claim of new complete resolution or a current-status certification. A known
answer reproduced for context does not become a new contribution. The catalogue's
source-status label is preserved; this attempt status does not assert that a
source-labelled solved question is open.
\end{document}
